
\begin{frame}[fragile]{Vues}

Toute requête produit une relation.

\vfill

Nommer cette requête c'est nommer la relation résultat.

\vfill

\alert{Une vue est une requête nommée}.

\vfill

On peut la considérer comme une relation \alert{calculée}
et l'interroger comme les autres. Très utile pour\,:

\vfill

\begin{redbullet}
   \item Filter une partie de la base
   \item Simplifier (et contrôler) les accès 
\end{redbullet}
\end{frame}


%%%%%%%%%%%%%%%%%%%%
\begin{frame}[fragile]{Création et interrogation d'une vue}

Exemple\,: on veut créer une ``vue''  de la base par région.

\vfill

\begin{minted}[baselinestretch=.9,
               fontsize=\small,
               framesep=2mm]{sql}
create view LogementCorse as 
        select *
        from Logement 
        where lieu='Corse'
\end{minted}

\vfill

On peut interroger la vue comme n'importe quelle table.

\vfill

\begin{minted}[baselinestretch=.9,
               fontsize=\small,
               framesep=2mm]{sql}
select * from LogementCorse
\end{minted}

\vfill

Le résultat de la requête
est réévalué à chaque fois que l'on accède à la vue. 

\end{frame}


%%%%%%%%%%%%%%%%%%%%
\begin{frame}[fragile]{Vue = ``macro'' pour des requêtes complexes}

On peut nommer des requêtes complexes souvent utilisées

\vfill

\begin{minted}[baselinestretch=.9,
               fontsize=\small,
               framesep=2mm]{sql}
create view VoyageursEnCorse as
select V.nom as nomVoyageur, L.nom as nomLogement
from LogementCorse as L, Séjour as S, Voyageur as V
where code = S.codeLogement
and   S.idVoyageur = V.idVoyageur
\end{minted}

\vfill

Notez\,: on construit des vues sur d'autres vues. Ici, on a tous
les clients des logements corses.
\end{frame}


%%%%%%%%%%%%%%%%%%%%
\begin{frame}[fragile]{Insertion dans une vue\,?}

Beaucoup de restrictions.


\vfill
\begin{redbullet}
  \item la vue doit être basée sur une seule table;
   \item toute colonne non référencée
    dans la vue doit pouvoir être mise à 
    \sqlkw{null} ou disposer d'une valeur par défaut;
  \item on ne peut pas mettre à jour un attribut
    qui résulte d'un calcul ou d'une opération.
\end{redbullet}

\vfill

On pourrait insérer dans \texttt{LogementCorse} mais pas dans \texttt{VoyageursEnCorse}

\end{frame}

%%%%%%%%%%%%%
\begin{frame}{À retenir}
Les \alert{vues} sont des requêtes nommées que l'on peut traiter
comme des tables.

\vfill

Elles permettent de restructurer ``en intention''  une base.

\begin{redbullet}
   \item  Pour faciliter l'accès (jointures pré-définies)
   \item Pour restreindre la visibilité des données (on ne donne accès qu'à la vue)
   \item Les insertions dans les vues sont soumises à fortes restrictions .
\end{redbullet}
\vfill

Un mécanisme très utile pour ``dénormaliser'' une base sans en subir les inconvénients. À suivre\,!

\end{frame}

